\chapter{Matrices and Systems}

\section{Systems of Equations as Matrices}\label{chap:syseqma}

\subsection{Systems of Equations, REF, and RREF}

We have talked a lot about matrices and vectors. Now we can use matrices and vectors to solve systems of equations. A system of linear equations can be written compactly using matrices and vectors. This matrix representation allows us to use algebraic tools such as row operations and matrix multiplication to analyze and solve the system more efficiently. Let us look at a very simple example from a previous section.

\[\begin{cases}
1x_1 + 2x_2 = -1 \\
2x_1 + 0x_2 = 4
\end{cases}\]

We can write this in the form of a matrix! This may come as a surprise to some of you. How can we do this? Let us write the coefficients as a matrix, the $x_n$'s as a column vector, and the constants as a column vector called $\vec{b}$, the right-hand-side vector.
\index{system of equations}
Here is the \index{coefficient matrix}coefficient matrix which we will call $A$:
\[A = \vecb{1 & 2 \\ 2 & 0}\]

Then we can write the column vector of $x$ values, which we will call $\vec{x}$
\[\vec{x} = \vecb{x_1 \\ x_2}\] 

And we can write the $\vec{b}$ as the following column vector:

\[
\vec{b} = \vecb{-1 \\ 4}
\]

Now, if we multiply $A$ and $\vec{x}$,
$$A\vec{x} = \vecb{1 & 2 \\ 2 & 0}\vecb{x_1 \\ x_2} = \vecb{1x_1 + 2x_2 \\ 2x_1 + 0x_2}$$

We get exactly the left-hand sides of our original system of equations!
Therefore, when we set this equal to the right-hand-side vector $\vec{b}$:
$$\vec{b}=\vecb{-1 \\ 4}$$
we get the \index{matrix equation}fundamental system of equation formula represented as a matrix Equation
\begin{equation}
       A\vec{x} = \vec{b}
       \label{eq:matrixform}
\end{equation}

This is a very important equation for linear algebra as a whole. Notice that the sizes of the matrices must follow the rules of matrix multiplication:
\[\underset{m \times n}{A} \cdot \underset{n \times 1}{\vec{x}} = \underset{m\times 1}{\vec{b}}\]

\begin{Exercise}[title=Matrix Equation, label=matrixeq1]
       Rewrite the systems of equations in matrix form and identify $A$, $\vec{x}$, and $\vec{b}$. 
       \[\begin{cases}
       -4x_1 + 9x_2 - 8x_3= 5\\
       -1x_1 + 0x_2 + 6x_3= 7
       \end{cases}\]
\end{Exercise}

\begin{Answer}[ref=matrixeq1]
     \[A = \vecb{-4 & 9 & -8 \\ -1 & 0 & 6}, \quad \vec{x} = \vecb{x_1 \\ x_2 \\ x_3}, \quad \vec{b} = \vecb{5 \\ 7}\]
     \[A\vec{x}=\vec{b} \implies \vecb{-4 & 9 & -8 \\ -1 & 0 & 6}\vecb{x_1 \\ x_2 \\ x_3} = \vecb{5 \\ 7}\]
\end{Answer}

\subsection{Row-Echelon Form}
This brings up the question, how can we solve for the $\vec{x}$? There are multiple ways, but most commonly we form an augmented matrix and solve for Row-Echelon Form. \index{row-echelon form}
\subsubsection{The Augmented Matrix}
\index{augmented matrix}
What is the augmented matrix? It is a way of representing the $A$ matrix and $\vec{b}$ vector. We write the $A$ matrix, separated by a vertical line from the $\vec{b}$ vector:
$$\begin{array}{c} [A \mid \vec{b}] =
\begin{bmatrix}
a_{11} & a_{12} & \cdots & a_{1n} & \vert & b_1 \\
a_{21} & a_{22} & \cdots & a_{2n} & \vert & b_2 \\
\vdots & \vdots &        & \vdots & \vert & \vdots \\
a_{m1} & a_{m2} & \cdots & a_{mn} & \vert & b_m
\end{bmatrix} \end{array}$$

The left block contains all the coefficients of the system of equations, and the right block contains only the solutions. Then, to solve for $\vec{x}$, we perform \textit{Elementary Row Operations} to find the \textit{Row-Echelon Form}, mathematically stated as $\text{rref}(A\vert \vec{b})$. Let us continue with the matrix vector set we were experimenting with:
$$A\vec{x} = \vecb{1 & 2 \\ 2 & 0}\vecb{x_1 \\ x_2} = \vecb{1x_1 + 2x_2 \\ 2x_1 + 0x_2}$$
This can be written as an augmented matrix:

$$\vecb{A\vert \vec{b}} = \vecb{1 & 2 &\vert & -1 \\ 2 & 0 & \vert & 4}$$


\subsubsection{Elementary Row Operations}
The Row-Echelon Form of a matrix is an \emph{equivalent} form of a matrix obtained through Gaussian and Gauss-Jordan elimination. There are lots of videos on this process online, so multiple videos included in your digital resources to walk you through the process. The rules of Gaussian Elimination are simple, only 3 different elementary row operations, or ERO's can be performed:

\begin{itemize}
       \item Swapping two rows
       \item Multiply a row by a nonzero scalar
       \item Add a multiple of one row to another row
\end{itemize}
Why do these operations preserve the given linear system? Swapping two rows only changes the order in which the equations are written and therefore has no effect on the solution set. Multiplying a row by a nonzero scalar corresponds to multiplying both sides of an equation by the same scalar, which does not alter its solutions. Finally, if two equations are satisfied by the same solution, then any linear combination of those equations is also satisfied by that solution; in particular, adding a multiple of one equation to another preserves the solution set. Before and after each operation, the augmented matrix contains exactly the same information as the system of equations and therefore represents the same solution set.

With these ERO's, we are trying to achieve a matrix with the following properties:\index{EROs|see{elementary row operations}}\index{elementary row operations}
\begin{itemize}
\item Leading entries (not necessarily 1) move to the right as you go down
\item Rows of zeros, if exist, are at the bottom
\item Each pivot has zeros in all positions below it in the same column.
\end{itemize}

Lets walk through the steps to simplify $\text{ref}(A\vert \vec{b})$:

\begin{align*}
\vecb{1 & 2 &\vert & -1 \\ 2 & 0 & \vert & 4} \\
R_2 \rightarrow R_2 - 2R_1 &\implies \vecb{1 & 2 &\vert & -1 \\ 0 & -4 & \vert & 6} \\
\end{align*}

As you can see, we set Row 2 to be Row 2 minus 2 times Row 1. This gives us a zero in the first column of Row 2. Technically, this matrix is now in Row-Echelon Form (REF). However, we can continue applying elementary row operations to obtain an even more structured form called the \emph{Reduced Row-Echelon Form}.

\subsubsection{Reduced Row Echelon Form}\index{reduced row-echelon form}
A matrix is in \emph{Reduced Row-Echelon Form} (RREF) if it satisfies all of the conditions for Row-Echelon Form, and in addition:

\begin{itemize}
\item Each leading entry is equal to $1$
\item Each leading entry is the only nonzero entry in its column
\end{itemize}

In other words, not only do pivots move to the right as we move down the rows, but each pivot column contains zeros both above and below the pivot, and each pivot is normalized to $1$. These additional conditions make RREF especially useful, since the solutions to the system can often be read directly from the matrix.

Continuing our example, we start with the matrix in REF:
\begin{align*}
\vecb{1 & 2 &\vert & -1 \\ 0 & -4 & \vert & 6}
\end{align*}

First, make the leading entry in Row 2 equal to $1$ by dividing the row by $-4$:
\begin{align*}
R_2 \rightarrow -\frac{1}{4}R_2 
&\implies 
\vecb{1 & 2 &\vert & -1 \\ 0 & 1 & \vert & -\frac{3}{2}}
\end{align*}

Next, eliminate the entry above the pivot in column 2:
\begin{align*}
R_1 \rightarrow R_1 - 2R_2 
&\implies 
\vecb{1 & 0 &\vert & 2 \\ 0 & 1 & \vert & -\frac{3}{2}}
\end{align*}

The matrix is now in Reduced Row-Echelon Form.

\subsubsection{Writing the Solution Vector}

Recall that the augmented matrix represents a system of linear equations. After reducing the matrix to Reduced Row-Echelon Form, we obtain
\begin{align*}
\vecb{1 & 0 &\vert & 2 \\ 0 & 1 & \vert & -\frac{3}{2}}
\end{align*}

Each row now corresponds directly to an equation in the variables $x_1$ and $x_2$:
\begin{align*}
x_1 &= 2 \\
x_2 &= -\frac{3}{2}
\end{align*}

Thus, the solution to the system can be written as the vector
\[
\vec{x}
=
\vecb{2 \\-\frac{3}{2}}
\]

Since every variable corresponds to a pivot column, the system has a \index{unique solution}\emph{unique solution}.

We have successfully solved the system of equations using matrix methods:
$$\vecb{1 & 2 \\ 2 & 0}\vecb{2 \\ -\frac{3}{2}} = \vecb{-1 \\ 4}$$

This is a fundamental technique in linear algebra, and you will use it often in this module and beyond!

\subsection{Larger RREF Example}

Let us solve a larger system by finding $\vec{x}$ in the matrix equation
\[
A\vec{x}=\vec{b}.
\]

\subsubsection{Write the System and Identify $A$, $\vec{x}$, and $\vec{b}$}

Consider the following system of equations:
\[
\begin{cases}
x_1 + x_2 + x_3 + x_4 = 4 \\
2x_1 + x_2 + 3x_3 + x_4 = 9 \\
x_1 + 3x_2 + 2x_3 + 2x_4 = 10 \\
3x_1 + x_2 + 2x_3 + 4x_4 = 13
\end{cases}
\]

The coefficient matrix, variable vector, and constant vector are:
\[
A=\vecb{
1 & 1 & 1 & 1\\
2 & 1 & 3 & 1\\
1 & 3 & 2 & 2\\
3 & 1 & 2 & 4
},
\qquad
\vec{x}=\vecb{x_1\\x_2\\x_3\\x_4},
\qquad
\vec{b}=\vecb{4\\9\\10\\13}.
\]

So the matrix equation is:
\[
A\vec{x}=\vec{b}.
\]

\subsubsection{Form the Augmented Matrix}

\[
\vecb{A\vert \vec{b}}
=
\vecb{
1 & 1 & 1 & 1 & \vert & 4\\
2 & 1 & 3 & 1 & \vert & 9\\
1 & 3 & 2 & 2 & \vert & 10\\
3 & 1 & 2 & 4 & \vert & 13
}.
\]

\subsubsection{Row-Reduce to REF}

Use the pivot in Row 1 to clear entries below it in column 1:
\begin{align*}
R_2 &\rightarrow R_2 - 2R_1 \\
R_3 &\rightarrow R_3 - R_1 \\
R_4 &\rightarrow R_4 - 3R_1
\end{align*}

\begin{align*}
\implies
\vecb{
1 & 1 & 1 & 1 & \vert & 4\\
0 & -1 & 1 & -1 & \vert & 1\\
0 & 2 & 1 & 1 & \vert & 6\\
0 & -2 & -1 & 1 & \vert & 1
}.
\end{align*}

Now use the pivot in Row 2 (column 2) to clear entries below it in column 2:
\begin{align*}
R_3 &\rightarrow R_3 + 2R_2 \\
R_4 &\rightarrow R_4 - 2R_2
\end{align*}

\begin{align*}
\implies
\vecb{
1 & 1 & 1 & 1 & \vert & 4\\
0 & -1 & 1 & -1 & \vert & 1\\
0 & 0 & 3 & -1 & \vert & 8\\
0 & 0 & -3 & 3 & \vert & -1
}.
\end{align*}

Next use the pivot in Row 3 (column 3) to clear the entry below it in column 3:
\begin{align*}
R_4 \rightarrow R_4 + R_3
\end{align*}

\begin{align*}
\implies
\vecb{
1 & 1 & 1 & 1 & \vert & 4\\
0 & -1 & 1 & -1 & \vert & 1\\
0 & 0 & 3 & -1 & \vert & 8\\
0 & 0 & 0 & 2 & \vert & 7
}.
\end{align*}

At this point, the matrix is in Row-Echelon Form (REF).

\subsubsection{Continue to RREF}

Make the pivot in Row 4 equal to $1$:
\begin{align*}
R_4 \rightarrow \frac{1}{2}R_4
\implies
\vecb{
1 & 1 & 1 & 1 & \vert & 4\\
0 & -1 & 1 & -1 & \vert & 1\\
0 & 0 & 3 & -1 & \vert & 8\\
0 & 0 & 0 & 1 & \vert & \frac{7}{2}
}.
\end{align*}

Clear the entries above the pivot in column 4:
\begin{align*}
R_1 &\rightarrow R_1 - R_4\\
R_2 &\rightarrow R_2 + R_4\\
R_3 &\rightarrow R_3 + R_4
\end{align*}

\begin{align*}
\implies
\vecb{
1 & 1 & 1 & 0 & \vert & \frac{1}{2}\\
0 & -1 & 1 & 0 & \vert & \frac{9}{2}\\
0 & 0 & 3 & 0 & \vert & \frac{23}{2}\\
0 & 0 & 0 & 1 & \vert & \frac{7}{2}
}.
\end{align*}

Make the pivot in Row 3 equal to $1$:
\begin{align*}
R_3 \rightarrow \frac{1}{3}R_3
\implies
\vecb{
1 & 1 & 1 & 0 & \vert & \frac{1}{2}\\
0 & -1 & 1 & 0 & \vert & \frac{9}{2}\\
0 & 0 & 1 & 0 & \vert & \frac{23}{6}\\
0 & 0 & 0 & 1 & \vert & \frac{7}{2}
}.
\end{align*}

Clear the entries above the pivot in column 3:
\begin{align*}
R_1 &\rightarrow R_1 - R_3\\
R_2 &\rightarrow R_2 - R_3
\end{align*}

\begin{align*}
\implies
\vecb{
1 & 1 & 0 & 0 & \vert & -\frac{10}{3}\\
0 & -1 & 0 & 0 & \vert & \frac{2}{3}\\
0 & 0 & 1 & 0 & \vert & \frac{23}{6}\\
0 & 0 & 0 & 1 & \vert & \frac{7}{2}
}.
\end{align*}

Make the pivot in Row 2 equal to $1$:
\begin{align*}
R_2 \rightarrow -R_2
\implies
\vecb{
1 & 1 & 0 & 0 & \vert & -\frac{10}{3}\\
0 & 1 & 0 & 0 & \vert & -\frac{2}{3}\\
0 & 0 & 1 & 0 & \vert & \frac{23}{6}\\
0 & 0 & 0 & 1 & \vert & \frac{7}{2}
}.
\end{align*}

Clear the entry above the pivot in column 2:
\begin{align*}
R_1 \rightarrow R_1 - R_2
\end{align*}

\begin{align*}
\implies
\vecb{
1 & 0 & 0 & 0 & \vert & -\frac{8}{3}\\
0 & 1 & 0 & 0 & \vert & -\frac{2}{3}\\
0 & 0 & 1 & 0 & \vert & \frac{23}{6}\\
0 & 0 & 0 & 1 & \vert & \frac{7}{2}
}.
\end{align*}

This matrix is now in Reduced Row-Echelon Form (RREF).

\subsubsection{Write the Solution Vector}

Reading each row as an equation gives:
\begin{align*}
x_1 &= -\frac{8}{3}\\
x_2 &= -\frac{2}{3}\\
x_3 &= \frac{23}{6}\\
x_4 &= \frac{7}{2}
\end{align*}

Thus the solution vector is:
\[
\vec{x}=\vecb{
-\frac{8}{3}\\[4pt]
-\frac{2}{3}\\[4pt]
\frac{23}{6}\\[4pt]
\frac{7}{2}
}.
\]

You can verify the solution by confirming that $A\vec{x}=\vec{b}$ using matrix multiplication.

Note that solution sets must be one of the following:
\begin{itemize}
       \item No solution (for example, two parallel lines that never meet),
       \item Exactly one unique solution (for example, two lines or planes that cross at one point),
       \item or Infinitely many solutions (for example, two equations describe the same line)
\end{itemize}
\subsection{Free Variables}

In the reduced row-echelon form (RREF) of a matrix, pivots determine which variables are constrained by the system.

\subsubsection{Zero rows and redundant equations}\index{zero row}
A row of all zeros in RREF represents a redundant equation. That row corresponds to the equation
\[
0x_1 + 0x_2 + \cdots + 0x_n = 0,
\]
which is always true and therefore imposes no restriction on the variables.  
Such a row arises when one of the original equations is a linear combination of the others.

Every row of zeros is a row without a pivot. Consequently, if a matrix has at least one zero row in RREF:
\begin{itemize}
    \item The number of pivots is less than the number of rows,
    \item The rank of the matrix is less than the number of rows,
    \item The rows of the matrix are linearly dependent.
\end{itemize}

\subsubsection{Free variables}
\index{pivot}Pivots are associated with \emph{columns}, not rows.  
A variable is called a \newterm{free variable}\index{free variable} if its column does not contain a pivot.

Thus:
\begin{itemize}
    \item A \emph{row} without a pivot (a zero row) indicates a redundant equation,
    \item A \emph{column} without a pivot corresponds to a free variable.
\end{itemize}

A free variable can take any value. When the system is consistent and has a free variable, it has infinitely many solutions.
\subsubsection{Infinitely Many Solutions}
\textbf{Example:}
Consider the augmented matrix in RREF:
\[
\vecb{1 & 0 & 0 & 8\\
      0 & 0 & 1 & -6}
\]

Labeling the variable columns (and ignoring the constants $c$ column), we write:
\[
\vecb{x_1 & x_2 & x_3 & c\\
      1   & 0   & 0   & 8\\
      0   & 0   & 1   & -6}
\]

The pivot positions are in the columns corresponding to \(x_1\) and \(x_3\).  
Since the column corresponding to \(x_2\) contains no pivot, \(x_2\) is a free variable.

The system of equations represented by this matrix is:
\[
\begin{aligned}
x_1 &= 8,\\
x_3 &= -6.
\end{aligned}
\]

The variable \(x_2\) does not appear in any equation and may take any real value.  
Since $x_1$ and $x_3$ are bound, we can say the solution is :
$$x_1 = 8,\qquad x_3 = -6,\qquad x_2 \in \mathbb{R}$$

written as a vector this is: $\vec{x} = \vecb{8 \\ x_2 \\ -6}$
or in set builder notation:
$\begin{array}{c} \left\{
\begin{pmatrix}
8\\
x_2\\
-6
\end{pmatrix}
:\; x_2 \in \mathbb{R}
\right\} \end{array}$

Thus, the presence of a column without a pivot leads to a free variable and infinitely many solutions.

\subsubsection{Inconsistent sets}
Let us look at another example:
\textbf{Example}
The augmented matrix
$$\vecb{
1&2&4&12\\
12&24&48&9\\
3&6&12&-4
}$$
simplifies to 
$$\vecb{
1&2&4&0\\
0&0&0&1\\
0&0&0&0}$$
we can interpret this As
$$\begin{aligned}
x_1 + 2x_2 + 4x_3 &= 0,\\
0x_1 + 0x_2 + 0x_3 &= 1.
\end{aligned}$$
Note that the equation $0=1$ is never true, so the system is \emph{inconsistent}\index{inconsistent set}, and there is \emph{no solution}.\index{no solution}

A row that is all zeros except for a nonzero entry in the last column signifies an inconsistent system.


\subsection{Basis, Dimension, Rank}
Let us define a few terms. These definitions are very basic, and will come back in the future:

\begin{mdframed}[frametitle = {Rank}, style = important]
\textbf{Rank} is the number of pivots in the RREF state of a matrix of size $m \times n$.
In other words, the number of nonzero rows in any row-echelon form of a matrix $A$ is denoted $\operatorname{rank}(A)$.
This has the constraints of:
\begin{itemize}
       \item $\operatorname{rank}(A) \le m$
       \item $\operatorname{rank}(A) \le n$
\end{itemize}
\end{mdframed}\index{rank}

\begin{mdframed}[frametitle = {Dimension}, style = important]
The \textbf{dimension} of the solution set of a system is the number of free variables. 
We calculate this with
$$\texttt{dimension} = \texttt{number of variables}-\texttt{rank}$$
or equivalently, for matrix $A$ of size $m \times n$
$$\texttt{dimension} = n - \operatorname{rank}(A)$$
\end{mdframed}\index{dimension}

\begin{mdframed}[frametitle = {Basis}, style = important]
A \textbf{basis} is a minimal set of direction vectors that describes all solutions.
Specifically, this consists of $n - \operatorname{rank}(A)$ vectors, one for each free variable
\end{mdframed}\index{basis}

In summary,
$$\begin{array}{lcl}
\text{rank} &=& \text{number of pivots} \\
\text{number of free variables} &=& n - \text{rank} \\
\text{dimension of solution set} &=& n - \text{rank} \\
\text{number of basis vectors} &=& n - \text{rank}
\end{array}$$

\subsection*{Exercises} % added this as a separator, maybe not needed
\begin{Exercise}[title=Augmented Matrix 1, label=infinitelymany]
Simplify the matrix:
$$\vecb{
1&2&-1&3\\
2&4&-2&6\\
1&1&0&2}$$
to RREF state and determine the type of solution set.
\end{Exercise}
\begin{Answer}[ref=infinitelymany]
Notice immediately that $R_2 = 2R_1$, which automatically means redundancy. 
$$\begin{bmatrix}
1&2&-1&3\\
2&4&-2&6\\
1&1&0&2
\end{bmatrix}
\;\xrightarrow{\substack{R_2\leftarrow R_2-2R_1\\ R_3\leftarrow R_3-R_1}}\;
\begin{bmatrix}
1&2&-1&3\\
0&0&0&0\\
0&-1&1&-1
\end{bmatrix}$$
$$\begin{array}{c} \xrightarrow{R_2\leftrightarrow R_3}\;
\begin{bmatrix}
1&2&-1&3\\
0&-1&1&-1\\
0&0&0&0
\end{bmatrix}
\;\xrightarrow{R_2\leftarrow -R_2}\;
\begin{bmatrix}
1&2&-1&3\\
0&1&-1&1\\
0&0&0&0
\end{bmatrix} \end{array}$$
$$\begin{array}{c} \xrightarrow{R_1\leftarrow R_1-2R_2}\;
\boxed{
\begin{bmatrix}
1&0&1&1\\
0&1&-1&1\\
0&0&0&0
\end{bmatrix}} \end{array}$$
The RREF is:
$$\vecb{
1&0&1&1\\
0&1&-1&1\\
0&0&0&0}$$
This gives us the system:
$$\begin{aligned}
x_1 + x_3 &= 1,\\
x_2 - x_3 &= 1.
\end{aligned}$$
So, 
$$x_1 = 1 - x_3,\qquad x_2 = 1 + x_3$$
making $x_3$ free.
We can write this as:
$\vec{x}
=
\vecb{
1\\1\\0
}
+
x_3
\vecb{
-1\\1\\1
},
\qquad x_3\in\mathbb{R}$
Equivalently, as a set:
$$\begin{array}{c} \left\{
\begin{pmatrix}
1-x_3\\
1+x_3\\
x_3
\end{pmatrix}
:\; x_3\in\mathbb{R}
\right\} \end{array}$$
Since $x_3$ can be anything, there are infinitely many solutions. The row of zeros shows that the system has a free variable, which is another way to see this.
\end{Answer}
\begin{Exercise}[title=Augmented Matrices 2, label=uniquesol]
Put the augmented matrix
$$\begin{bmatrix}
1&-1&2&0\\
0&1&-1&3\\
2&1&0&1
\end{bmatrix}$$
into RREF and determine its solution.
\end{Exercise}
\begin{Answer}[ref=uniquesol]
Start with $R_3 \leftarrow R_3 - 2R_1$ to clear the first entry in row 3:
$$\begin{bmatrix}
1&-1&2&0\\
0&1&-1&3\\
0&3&-4&1
\end{bmatrix}$$
Next, $R_3 \leftarrow R_3 - 3R_2$ gives the row echelon form:
$$\begin{bmatrix}
1&-1&2&0\\
0&1&-1&3\\
0&0&-1&-8
\end{bmatrix}$$
Then $R_3 \leftarrow -R_3$, $R_2 \leftarrow R_2 + R_3$, and $R_1 \leftarrow R_1 - 2R_3$ give
$$\begin{bmatrix}
1&-1&0&-16\\
0&1&0&11\\
0&0&1&8
\end{bmatrix}$$
Finally, $R_1 \leftarrow R_1 + R_2$ gives the RREF:
$$\begin{bmatrix}
1&0&0&-5\\
0&1&0&11\\
0&0&1&8
\end{bmatrix}$$
This gives a unique solution:
$$x_1=-5,\quad x_2=11,\quad x_3=8$$
or as a vector
\(
\vec{x} = \vecb{-5\\11\\8}
\)
\end{Answer}

\begin{Exercise}[title=Augmented Matrices 3, label=nosol]
Determine the solution for the following set:
$$\begin{bmatrix}
1&2&3&4\\
2&4&6&9\\
0&0&0&1
\end{bmatrix}$$
\end{Exercise}
\begin{Answer}[ref=nosol]
Notice that one row of the matrix states $\vecb{0&0&0&1}$, which implies $0x_1 + 0x_2+0x_3=1$, which is not possible. Therefore the matrix represents an inconsistent system and there is no solution.
\end{Answer}
\begin{Exercise}[title=Rank 1, label=rank1]
Find the rank of
$$\begin{bmatrix}
1&2&3\\
2&4&6\\
3&6&9
\end{bmatrix}$$
\end{Exercise}
\begin{Answer}[ref=rank1]
The RREF of this matrix is
$$\begin{bmatrix}
1&2&3\\
2&4&6\\
3&6&9
\end{bmatrix}
\xrightarrow{\substack{R_2\leftarrow R_2-2R_1\\ R_3\leftarrow R_3-3R_1}}
\begin{bmatrix}
1&2&3\\
0&0&0\\
0&0&0
\end{bmatrix}$$
Note that there are two rows of zeroes, and the first column has a 1. This implies two dependent rows ($R_2 = 2 R_1 \qquad R_3 = 3 R_1$), and with only 1 pivot column, $\operatorname{rank}(A)=1$. 
\end{Answer}

\begin{Exercise}[title=Rank 2,label=rank2]
Let $A$ be a $4\times 6$ matrix with exactly $3$ pivots in $\operatorname{RREF}$.
\begin{itemize}
       \item What is $\operatorname{rank}(A)$?
       \item How many free variables are there?
\end{itemize}
\end{Exercise}
\begin{Answer}[ref=rank2]
By definition in this chapter:
$$\operatorname{rank}(A)=\text{number of pivots}=3$$

A $4 \times 6$ matrix has $n=6$ columns, so 
$$\text{\# free variables} = n-\operatorname{rank}(A)=6-3=3.$$
So, $\text{free variables} = 3$
\end{Answer}

\begin{Exercise}[title=Free Variables,label=free1]
The RREF of a system has the form:
$$\begin{bmatrix}
1&0&-2&0&3\\
0&1&1&0&-1\\
0&0&0&1&4\\
0&0&0&0&0
\end{bmatrix}$$
\begin{enumerate}
       \item Identify the pivot variables
       \item Identify the free variables
       \item State the dimension of the solution set
\end{enumerate}
\end{Exercise}
\begin{Answer}[ref=free1]
There are $n=4$ variables $(x_1,x_2,x_3,x_4)$. 
\begin{enumerate}
       \item Pivot columns are columns $1,2,4$, so the pivot columns are $x_1,\ x_2,\ x_4$. 
       \item Free variables: the only non-pivot column is column $3$, so the free variable is $x_3$
       \item The dimension of the solution set is the number of free variables: $1$. Equivalently, $\text{dimension}=n-\text{rank}=4-3=1$
\end{enumerate}
\end{Answer}
\subsection{Creating Matrices in Python}
Recall that matrices can be represented by 2D arrays in Python. Let us explore this idea further with the module \mintinline{python}|numpy|. Create a file called \filename{matrices\_creation.py} and enter this code:
\begin{minted}{python}
# import the python module that supports matrices
import numpy as np
# Use the function np.array to define a matrix that 
# contains specific values that you supply.
A = np.array([[ 5, 1, 3], 
              [ 1, -1, 8], 
              [ 6, 2, 1]])
# The transpose function returns 
A.transpose()
\end{minted}
When you run it, you should see:
\begin{minted}{python}
array([[ 5, 1, 6], 
       [ 1, -1, 2], 
       [ 3, 8, 1]])
\end{minted}
As you can see, $A\neq A^T$, so A is not symmetric.
Try another: 
\begin{minted}{python}
# create a matrix, B
B = np.array([[ 5, 1, 6], 
              [ 1, -1, 2], 
              [ 6, 2, 1]])
B.transpose()
\end{minted}
When you run it, you should see:
\begin{minted}{python}
array([[ 5, 1, 6], 
       [ 1, -1, 2], 
       [ 6, 2, 1]])
\end{minted}
B is symmetric. You can actually transpose any matrix using this function, but a matrix cannot be symmetric unless it is square. 

Try this code to see what happens when you transpose a rectangular matrix. 
\begin{minted}{python}
# create a matrix, J
J = np.array([[ 5, 1, 3, 0], 
              [ 1, -1, 8, 11], 
              [ 6, 2, 1,-7]])
J.transpose()
\end{minted}
Note that transposing a rectangular matrix changes its dimension from 3 by 4 to 4 by 3. You should see a transposed matrix, but it is not symmetric.
\begin{minted}{python}
array([[ 5,  1,  6],
       [ 1, -1,  2],
       [ 3,  8,  1],
       [ 0, 11, -7]])
\end{minted}

\subsubsection{Creating Special Matrices in Python}
 Use the same file to add this code for creating a zero matrix.
\begin{minted}{python}
# create an 8 by 10 Zero matrix.
 C = np.zeros((8, 10))
 C
\end{minted}
When you run it, you should see:
\begin{minted}{python}
array([[0., 0., 0., 0., 0., 0., 0., 0., 0., 0.],
       [0., 0., 0., 0., 0., 0., 0., 0., 0., 0.],
       [0., 0., 0., 0., 0., 0., 0., 0., 0., 0.],
       [0., 0., 0., 0., 0., 0., 0., 0., 0., 0.],
       [0., 0., 0., 0., 0., 0., 0., 0., 0., 0.],
       [0., 0., 0., 0., 0., 0., 0., 0., 0., 0.],
       [0., 0., 0., 0., 0., 0., 0., 0., 0., 0.],
       [0., 0., 0., 0., 0., 0., 0., 0., 0., 0.]])
\end{minted}
Add the following code to create an 8 by 8 Identity matrix.
\begin{minted}{python}
# create an 8 by 8 Identity matrix 
 D = np.eye(8)
 D
\end{minted}
When you run it, you should see:
\begin{minted}{python}
array([[1., 0., 0., 0., 0., 0., 0., 0.],
       [0., 1., 0., 0., 0., 0., 0., 0.],
       [0., 0., 1., 0., 0., 0., 0., 0.],
       [0., 0., 0., 1., 0., 0., 0., 0.],
       [0., 0., 0., 0., 1., 0., 0., 0.],
       [0., 0., 0., 0., 0., 1., 0., 0.],
       [0., 0., 0., 0., 0., 0., 1., 0.],
       [0., 0., 0., 0., 0., 0., 0., 1.]])
\end{minted}
As you progress in your studies, you will learn the importance of diagonal matrices and of extracting the diagonal of a matrix. Let us see how to extract a diagonal, then create a diagonal matrix.
\begin{minted}{python}
# create a  matrix 
W = np.array([[1, 2, 3, 4], 
              [5, 6, 7, 8], 
              [-8, -7, -6, -5],
              [-4, -3, -2, -1]])
\end{minted}
Extract the main diagonal using \index{numpy!diag}np.diag(<array>,<diagonal to extract>). 
Passing 0 as the second parameter specifies the main diagonal. A positive value extracts a diagonal from the upper part. A negative value extracts a diagonal from the lower part. Run this code then experiment passing other values to see what you get.
\begin{minted}{python}
print(np.diag(W,0))
\end{minted}
When you run it, you should see:
\begin{minted}{python}
array([ 1,  6, -6, -1])
\end{minted}
You can also use np.diag() to create a diagonal matrix from a 1D array. In this case, do not pass a second parameter.
\begin{minted}{python}
Q = np.array([1, 2, 3])
DiagArray = np.diag(Q)
print(DiagArray)
\end{minted}
When you run it you should see;
\begin{minted}{python}
[[1 0 0]
 [0 2 0]
 [0 0 3]]
\end{minted}
Python has functions for extracting upper and lower triangular matrices. Try these:
\begin{minted}{python}
print(np.triu(W))
print(np.tril(W))
\end{minted}
You should see:
\begin{minted}{python}
[[ 1  2  3  4]
 [ 0  6  7  8]
 [ 0  0 -6 -5]
 [ 0  0  0 -1]]
 
[[ 1  0  0  0]
 [ 5  6  0  0]
 [-8 -7 -6  0]
 [-4 -3 -2 -1]]
\end{minted}

\subsection{Conclusion}
In this module, we have explored the fundamental concepts of matrices, including their types, properties, and how they can be used to represent and solve systems of linear equations. We have also learned how to perform row operations to convert matrices into Row-Echelon Form (REF) and Reduced Row-Echelon Form (RREF), which are essential techniques in linear algebra. In the next chapter~\ref{chap:subspaces}, we will talk deeper about the concepts of vector spaces and subspaces and how to find RREF, building on our understanding of matrices and their applications.


\section{Applications of Matrices in System Form}

In the chapter on linear combinations, we saw that we can linearly combine vectors to create other vectors. Consider 3 vectors:

$$\textbf{x} = \begin{bmatrix}
-1\\
2\\
0
\end{bmatrix}\quad \textbf{y} = \begin{bmatrix}
-1\\
2\\
1
\end{bmatrix}\quad \textbf{z} = \begin{bmatrix}
-2\\
1\\
0
\end{bmatrix}$$

We can write a linear combination of these vectors:
$$c\textbf{x} + d\textbf{y} + e\textbf{z}$$

Which we can expand to show the vectors:
$$c \begin{bmatrix}
-1\\
2\\
0
\end{bmatrix} + d \begin{bmatrix}
-1\\
2\\
1
\end{bmatrix} + e \begin{bmatrix}
-2\\
1\\
0
\end{bmatrix} = \begin{bmatrix}
-c - d - 2e\\
2c + 2d + e\\
d
\end{bmatrix}$$

We can also represent this combination with a \index{linear combination!matrix form}matrix where each column is one of the vectors:
$$\begin{bmatrix}
-1 & -1 & -2\\
2 & 2 & 1\\
0 & 1 & 0
\end{bmatrix} \cdot \begin{bmatrix}
c\\
d\\
e
\end{bmatrix} = \begin{bmatrix}
-c - d - 2e\\
2c + 2d + e\\
d
\end{bmatrix}$$

%\subsection{Thinking about Rows}
	
%\subsection{Thinking about Columns}


%\section{Systems of Linear Equations}
%matrix-vector multiplication, representing linear combination
%system of linear equations connection
%linear combo example
%system of equations example (trail mix, be more explicit about eqs and representation)
%practice
%	use matrix multiplication to find...
%	grading question moved here

\subsection{Trail Mix for Mars}\index{trail mix example}
Let us look at an applied problem. Three astronauts, Pat, Kai, and River, are getting ready for a trip to Mars. NASA food service is preparing trail mix for the voyage, tailored to each astronaut's taste. The chef needs to submit a budget based on the cost of the trail mix for each astronaut. The mix is made up of raisins, almonds, and chocolate.

Pat prefers a raisins:almonds:chocolate ratio of 6:10:4, Kai likes 2:3:15, and River wants 14:1:5. The chef can buy a kg of raisins for \$7.50, a kg of almonds for \$14.75, and a kg of chocolate for \$22.25. Assuming each astronaut will get 20 kg of trail mix, which astronaut will cost more to feed?

First, set up a matrix to represent the raisins:almonds:chocolate ratios. (Conveniently, these ratios already add to 20.)

$$MixRatios = \begin{bmatrix}
6 & 10 & 4  \\
2 & 3 & 15 \\
14 & 1 & 5
\end{bmatrix}$$

Use a vector to represent the cost of each item:
$$IngredientCost = 
\begin{bmatrix}
7.50 \\
14.75 \\
22.25
\end{bmatrix}$$

To find the cost of trail mix for each astronaut, we simply find the \index{dot product!application}dot product between the mix ratios and the ingredient costs to get:

Pat =  \$281.50 \newline
Kai =  \$393.00 \newline
River = \$231.00 \newline

So Kai will cost the most to feed.


\begin{Exercise}[title={Vector Matrix Multiplication}, label=vector-matrix-multiply01]
Multiply the array $A$ with the vector $v$. Compute this by hand, and make sure to show your computations. 
$$A = \begin{bmatrix}
1 & -2 & 3 & 5  \\
-4 & 2 & 7 & 1 \\
3 & 3  & -9 & 1
\end{bmatrix}$$

$$v = 
\begin{bmatrix}
 2 \\
 2 \\
 6 \\
 -1
 \end{bmatrix}$$
\end{Exercise}

\begin{Answer}[ref=vector-matrix-multiply01]
$$Av = \vecb{11 \\ 37 \\ -43}$$
\end{Answer}

\begin{Exercise}[title = {Using Vector Matrix Multiplication}, label = vmm02]
A college professor offers three different methods of determining a student's final grade. In method A, the student's grade is 20\% based on attendance, 50\% homework, 15\% midterm, and 15\% final. This professor knows many students can learn the material without attending every class, so with method B the student's grade is 50\% homework, 20\% midterm, and 30\% final. Last, the professor knows some students do not do the homework but still show they understand the material by doing well on the tests. With method C, a student's grade is 40\% midterm and 60\% final. \index{weighted grades}The professor uses whatever method results in the highest grade to determine each student's final grade. 

Suppose Suzy has attended 65\% of classes, has an average homework grade of 30\%, earned a 80\% on the midterm, and earned a 75\% on the final. What final grade will her professor post?
\end{Exercise}

\begin{Answer}[ref = vmm02]
The different methods can be represented in a matrix:
$$\begin{bmatrix}
0.20 & 0.50 & 0.15 & 0.15\\
0 & 0.50 & 0.20 & 0.30\\
0 & 0 & 0.4 & 0.6
\end{bmatrix}$$

And Suzy's individual grades can be represented by a vector:
$$\begin{bmatrix}
65\\
30\\
80\\
75
\end{bmatrix}
$$

To see the results of the three different methods, we can multiply the matrix and the vector:
$$\begin{bmatrix}
0.20 & 0.50 & 0.15 & 0.15\\
0 & 0.50 & 0.20 & 0.30\\
0 & 0 & 0.4 & 0.6
\end{bmatrix}
\cdot 
\begin{bmatrix}
65\\
30\\
80\\
75
\end{bmatrix} = 
\begin{bmatrix}
0.2(65)+0.5(30)+0.15(80) + 0.15(75)\\
0(65) + 0.5(30) + 0.2(80) + 0.3(75)\\
0(65) + 0(30) + 0.4(80) + 0.6(75)
\end{bmatrix}$$

Which yields:
$$\begin{bmatrix}
51.25\\
53.5\\
77
\end{bmatrix}$$

Since method C yields the highest grade, the professor will post a final grade of 77. 
\end{Answer}

\subsubsection{Vector-Matrix Multiplication in Python}\index{matrix-vector multiplication}
Most real-world problems use very large matrices, where it becomes impractical to do calculations by hand. As long as you understand how matrix-vector multiplication is performed, you will be equipped to use a computing language, like Python, to do the calculations for you. 

Create a file called \filename{vectors\_matrices.py} and enter this code:

\begin{minted}{python}
# import the python module that supports matrices
import numpy as np

# create an array
a = np.array([[5, 1, 3, -2], 
              [1, -1, 8, 4], 
              [6, 2, 1, 3]])

# create a vector 
b = np.array([1, 2, 3, -8])

# calculate the dot product
print(a.dot(b))
\end{minted}

When you run it, you should see:
\begin{minted}{python}
[32 -9 -11]
\end{minted}

\subsection{Where to Learn More}
Watch this video from Khan Academy about matrix-vector products: \url{https://rb.gy/frga5}
